\documentclass[11pt,oneside,openany]{book}
\usepackage[paperwidth=6.5in,paperheight=9.5in,inner=.72in,outer=.72in,top=.75in,bottom=.78in,headheight=14pt]{geometry}
\usepackage{fontspec}
\setmainfont{lmroman10-regular.otf}[BoldFont=lmroman10-bold.otf,ItalicFont=lmroman10-italic.otf,BoldItalicFont=lmroman10-bolditalic.otf,SmallCapsFont=lmromancaps10-regular.otf]
\usepackage[main=english]{babel}
\usepackage{amsmath,amssymb,graphicx,microtype,fancyhdr,needspace,tikz}
\usepackage[unicode,hidelinks]{hyperref}
\hypersetup{pdftitle={Carl Friedrich Gauss: Works, Volume V. Attraction, Mechanics, Fluid Equilibrium, Terrestrial Magnetism and General Theorems, Dioptrical Investigations, Notices of the Author’s Own Papers, Introduction to the Magnetic Association's 1836 Observations, A New Aid for Magnetic Observations, A New Instrument for Measuring Changes in Magnetic Intensity, Instructions for Determining the Duration of a Magnetic Needle’s Swing, On a Way to Make Deflections Easier to Observe, On Determining the Constants of the Bifilar Magnetometer, Instructions for Determining the Magnetic Effect Exerted by a Bar Magnet at a Distance, On Using the Magnetometer to Determine Absolute Declination, Observations of Magnetic Inclination in Göttingen, Fundamental Equations for the Motion of Heavy Bodies on the Rotating Earth, On Achromatic Doublet Objectives, Letter from Gauss to Brandes, Adjusting the Knife Edges of a Balance, Physical Observations, Northern Lights on 7 January 1831, Magnetic Observatory in Göttingen (August 1834), Observations of Magnetic Variation in Göttingen and Leipzig, Magnetic Observations (March 1835), Magnetic Needle in Copenhagen and Milan, Letter from Gauss to Schumacher (1836), Procedure for Scheduled Observations, Three Years of Daily Declination Observations, Plots and Numerical Readings for the Scheduled Observations of 1836, Plots and Numerical Readings for the Scheduled Observations of 1837, On Wilkes’s Observations near the Magnetic South Pole, Review of Benzenberg’s On Dalton’s Theory, Review of Fischer’s Instructions for Making Artificial Magnets, Review of the Magnetic Association’s 1836 Results, The Mathematical and Physical Institute in Marburg, On the Mathematical Theory of Electrodynamic Effects, §§ 1–21, Spherical Harmonics, On Using the Comparator, Effect of a Luminous Point, Schering’s Remarks, Historical Contents and Imprint},pdfauthor={Carl Friedrich Gauss}}
\setlength{\parindent}{1.2em}\setlength{\parskip}{0pt plus .3pt}
\setlength{\emergencystretch}{1.4em}
\pagestyle{fancy}\fancyhf{}
\fancyhead[L]{\small THEORY OF ATTRACTION}
\fancyfoot[C]{\thepage}\renewcommand{\headrulewidth}{.25pt}
\newcommand{\SourceLeaf}[1]{\hypertarget{native-#1}{}\label{native-#1}}
\newcommand{\Article}[1]{\par\medskip\Needspace{4\baselineskip}\pdfbookmark[1]{Art. #1}{attractionis-art-#1}\centerline{\textbf{#1.}}\nobreak\smallskip}
\renewcommand{\thefootnote}{*)}
\newcommand{\HistoricalNote}[1]{\footnote{\csname HistoricalNoteText#1\endcsname}}
\newcommand{\BareRoot}[2][]{\mathord{{}^{#1}\!\surd}#2}
\widowpenalty=10000\clubpenalty=10000\displaywidowpenalty=10000
\raggedbottom
